% Einheit 4 von 6 – Symmetrie (bank/funktionsklassen-und-eigenschaften/e4.jsonl, zusammenbau v0.1)
%% TODO Zeitmarke Einheit 4: Marken-Zeile nicht lesbar
%% TODO Zweigzeile Teil 1: was hier gelernt wird (Satz in Schülersprache aus dem Einheitstitel) – Einheit 4
\einheitenkopf[e4]{Einheit 4 von 6 · Symmetrie}
\zweigzeile{Abitur GK}
\setcounter{aufgabe}{16}

%% TODO Titel: Ich-kann-Satz fehlt in der Bank (Platzhalter: Kettenname) – Nr. 17
%% TODO Anweisung: ein Satz über der ersten Teilaufgabe fehlt in der Bank – Nr. 17
\begin{aufgabe}{Symmetrie}
%% TODO \rechenplatz{n}: Schrittzahl des Grundfalls fehlt in der Bank (nur bei fester Schreibform, 2.3 b)
\begin{teile}
\teil $f(x) = 2x^6 - x^2 + 5$ – gerade oder ungerade Exponenten? \\ \kreuz{achsensymmetrisch} \\ \kreuz{punktsymmetrisch} \\ \kreuz{keines von beiden}
\teil $f(x) = x^4 - 6x^2 + 2$ – Symmetrie? Begründe am Term.
\teil $f(x) = x^2 \cdot e^{-x^2}$ – Symmetrie? Begründe über $f(-x)$.
\teil $f_a(x) = a x^3 - x + a$ – Symmetrie von $f_0$?
\teil $f(x) = x^4 - 8x^2 + 7$ – Symmetrie mit Begründung und Schnittpunkt mit der y-Achse?
\teil $f$ ist achsensymmetrisch zur y-Achse, hat die Nullstelle $2{,}5$ und den Hochpunkt $H(1{,}5 | 4)$. Weitere Nullstelle und weiterer Hochpunkt?
\teil Ein symmetrischer Brückenbogen ist am Boden $24$ m breit und $6$ m hoch. Das Koordinatensystem wird so gelegt, dass die y-Achse die Symmetrieachse ist und der Boden auf der x-Achse liegt (1 Längeneinheit entspricht 1 m). Koordinaten der Fußpunkte und des höchsten Punkts?
\teil Die Gerade $y = c$ schneidet den Graphen von $f(x) = 0{,}5x^2$ in zwei Punkten, die $6$ Längeneinheiten auseinanderliegen. Wert von $c$? c = \leerfeld
\teil $f(x) = 0{,}5x^4 - 4{,}5x^2 + 2$. Begründe, warum der Graph symmetrisch zur y-Achse ist, und gib das Verhalten der Funktionswerte für $x \to -\infty$ und $x \to +\infty$ an \hfill (FHR 2026).
\end{teile}
\end{aufgabe}

%% TODO Titel: Ich-kann-Satz fehlt in der Bank (Platzhalter: Kettenname) – Nr. 18
%% TODO Anweisung: ein Satz über der ersten Teilaufgabe fehlt in der Bank – Nr. 18
\begin{aufgabe}{Symmetrie – Fehler finden, Begründen}
\begin{teile}
\teil Ich finde den Fehler: „$f(x) = x^4 - 5x^2 + 5$ ist nicht symmetrisch, weil die $5$ kein $x$ hat und deshalb ungerade zählt.“
\teil Begründe: Das Absolutglied zählt bei der Symmetrie als gerader Exponent.
\end{teile}
\end{aufgabe}
